\begin{abstract}
\noindent
We study the intraday response of \eurusd{} and the Nasdaq-100 index
to 2{,}449 unscheduled news events from a public real-time newsfeed
over a thirteen-month sample. Each event is independently labelled by
a large language model and benchmarked against the Loughran-McDonald
dictionary and FinBERT-tone. Using a matched, hour-of-day and
day-of-week stratified baseline with cluster-robust standard errors
and uniform Benjamini-Hochberg false discovery rate adjustment, we
find that news events are associated with intraday magnitude
$1.3\times$ to $1.9\times$ the baseline level on every asset-window
cell. The directional content of LLM sentiment is modest (49--54\%
hit rate) and likely too small to overcome typical intraday
transaction costs, but the LLM
outperforms the dictionary baseline by about 1--6 percentage points,
is competitive with FinBERT across primary windows, and produces
asset-specific sentiment on 73.5\% of events. On the
Financial PhraseBank external benchmark, the LLM attains 84.6\%
accuracy and Cohen's $\kappa = 0.72$ against gold human labels,
outperforming FinBERT by 5.4 percentage points. The pre-event
absolute return matches the post-event return in magnitude,
indicating that the public headline timestamp is not the
informational event time. The full pipeline is released with seeded
reproducibility.
\end{abstract}
